\subsubsection{Análisis de Alternativas y Justificación [Pendientes: 1. documentar comparación de simuladores y herramientas aún por elegir]}\label{sec:spmp-8-2-1}
WFDS se mantiene como la elección que el equipo considera más adecuada para el tipo de representación física y de escenarios que pretende preparar. La decisión vigente es utilizarlo de forma exclusiva, mientras que FARSITE se conserva como antecedente del proceso de formulación. Esta justificación debe vincularse con las necesidades y restricciones de las secciones \ref{sec:spmp-6-2} y \ref{sec:spmp-6-3}, y completarse con la evidencia técnica y las comparaciones ya disponibles, porque la preferencia del equipo no constituye una demostración de que sea el mejor simulador para todos los contextos. La ejecución del piloto será una comprobación de viabilidad de la elección y no una comparación retrospectiva ficticia con herramientas que no se hayan ensayado~\cite{diseno2026}.

UML se elige porque aporta convenciones compartidas para comunicar el diseño y porque sus conceptos son conocidos por el equipo como estudiantes de Ingeniería de Sistemas de la Pontificia Universidad Javeriana. Aldine se adopta como reemplazo de Overleaf para disponer de un espacio centralizado de documentos sin depender del pago de uso de esa plataforma, aprovechando el servidor disponible. Esta motivación no presupone que el computador, su operación o sus respaldos carezcan de costos. Git, por su parte, permitirá conservar la evolución del código de la aplicación y relacionar cambios con versiones comprobadas. Las decisiones anteriores ofrecen razones de selección, aunque la comparación de otras tecnologías de desarrollo, pruebas y almacenamiento todavía requiere fuentes y criterios documentados~\cite{uml,servidor2026,adaptacion2026}.
